Placing an atom in an external electric field splits its energy levels. The number of distinct levels thereby produced, and the symmetry of their states, can be found by the method described in §~96. To do this, the reducible $(2J+1)$-dimensional representation of the symmetry group of the external field, furnished by the functions $\psi_{JM}$, must be decomposed into irreducible representations of that group. We therefore need the characters of the representation furnished by the functions $\psi_{JM}$. Since irreducible-representation characters have the same value for all elements of a class, it is enough to consider rotations about a single axis, chosen as the $z$ axis. Under a rotation through $\varphi$ about $z$, the wavefunctions $\psi_{JM}$ are, as we know, multiplied by $e^{iM\varphi}$, where $M$ is the projection of the angular momentum on that axis. The transformation matrix for the functions $\psi_{JM}$ is consequently diagonal, and its character is
\[
 \chi^{(J)}(\varphi)=\sum_{M=-J}^{J}e^{iM\varphi}
 =\frac{e^{i(J+1)\varphi}-e^{-iJ\varphi}}{e^{i\varphi}-1},
\]
or\footnote{To avoid misunderstanding, we emphasize that this formula uses a parametrization of group elements different from the Euler-angle parametrization: a transformation is specified by the direction of the rotation axis and by the angle $\varphi$ about it. With this parametrization it can be shown, for example, that the integration in (98.2) is to be performed with the element $2(1-\cos\varphi)\,d\varphi\,do$, where $do$ is the element of solid angle for the direction of the rotation axis.}
\begin{equation}
 \chi^{(J)}(\varphi)=\frac{\sin\!\left((J+\tfrac12)\varphi\right)}{\sin(\varphi/2)}.
 \tag{98.3}
\end{equation}
Under inversion $I$, on the other hand, all the functions $\psi_{JM}$ with different $M$ behave alike: they are multiplied by $+1$ or by $-1$ according as the atomic state is even or odd.

